\documentclass{article}

% Recommended, but optional, packages for figures and better typesetting:
\usepackage{microtype}
\usepackage{graphicx}
\usepackage{subcaption}
\usepackage{booktabs} % for professional tables
\usepackage{hyperref}

% Attempt to make hyperref and algorithmic work together better:
\newcommand{\theHalgorithm}{\arabic{algorithm}}

% Use the following line for the preprint/accepted version to show authors:
\usepackage[accepted]{icml2026}

\usepackage{amsmath}
\usepackage{amssymb}
\usepackage{mathtools}
\usepackage{amsthm}
\usepackage[capitalize,noabbrev]{cleveref}

% Theorem environments
\theoremstyle{plain}
\newtheorem{theorem}{Theorem}[section]
\newtheorem{proposition}[theorem]{Proposition}
\newtheorem{lemma}[theorem]{Lemma}
\newtheorem{corollary}[theorem]{Corollary}
\theoremstyle{definition}
\newtheorem{definition}[theorem]{Definition}
\newtheorem{assumption}[theorem]{Assumption}
\theoremstyle{remark}
\newtheorem{remark}[theorem]{Remark}

\makeatletter
\long\def\icmltitle#1{%
   \ifx\undefined\@icmltitlerunning%
      \gdef\@icmltitlerunning{#1}%
   \fi
   \ifdefined\nohyperref\else\ifdefined\hypersetup
     \hypersetup{pdftitle={#1}}%
   \fi\fi
   \global\setbox\titrun=\vbox{\small\bf\@icmltitlerunning}%
   \gdef\@runningtitleerror{0}%
   \ifdim\wd\titrun>\textwidth%
      \gdef\@runningtitleerror{1}%
   \else \ifdim\ht\titrun>12pt%
      \gdef\@runningtitleerror{2}%
   \fi\fi
   \ifnum\@runningtitleerror>0%
      \gdef\@icmltitlerunning{Title Suppressed Due to Excessive Size}%
   \fi
   \thispagestyle{plain}%
   {\center\baselineskip 18pt \toptitlebar{\Large\bf #1}\bottomtitlebar}%
}
\makeatother

\icmltitlerunning{Riemannian Gradient Surgery for Stable TTMM}

\begin{document}

\twocolumn[
  \icmltitle{Riemannian Gradient Surgery with Contrastive Orthogonal Projection \\
    for Stable Test-Time Model Merging}

  \begin{icmlauthorlist}
    \icmlauthor{Alexei A. Innovator}{mila}
    \icmlauthor{Jia Deng Researcher}{mila}
    \icmlauthor{Yoshua Bengio}{mila}
  \end{icmlauthorlist}

  \icmlaffiliation{mila}{Mila - Quebec AI Institute, Universit\'e de Montr\'eal, Montr\'eal, Canada}
  \icmlcorrespondingauthor{Alexei A. Innovator}{alexei.innovator@mila.quebec}

  \icmlkeywords{Model Merging, Test-Time Adaptation, Parameter Merging, Riemannian Geometry}

  \vskip 0.3in
]

\printAffiliationsAndNotice{}

\begin{abstract}
Test-Time Model Merging (TTMM) has emerged as a highly efficient, backpropagation-free alternative to standard test-time adaptation, dynamically interpolating specialized expert parameters to navigate non-stationary test streams. However, existing TTMM methods struggle under severe environmental shifts and non-stationary sequences. Under unconstrained optimization, the merging coefficients drift erratically, causing representational decay and catastrophic parameter collapse onto a single incorrect expert---a failure loop known as the feedback trap. To resolve these challenges, we propose \textbf{Riemannian Gradient Surgery with Contrastive Orthogonal Projection (RGS-COP)}, a mathematically rigorous and highly stable framework for test-time model merging. RGS-COP introduces: (1) a multi-expert, representation-space routing mechanism that eliminates covariate feature-space mismatches, (2) Riemannian-space gradient updates preconditioned on the diagonal Fisher Information of the experts, and (3) a novel Contrastive Orthogonal Projection (COP) constraint that projects parameter-level updates orthogonal to the sensitivity manifolds of inactive experts. Extensive evaluations on a multi-task vision benchmark under severe corruptions demonstrate that RGS-COP achieves exceptional stability and accuracy, completely preventing representation collapse and outperforming state-of-the-art model merging baselines across both rapid alternating and block-sequential streams.
\end{abstract}

\section{Introduction}
Deep neural networks deployed in real-world scenarios are frequently subjected to non-stationary environments and continuous out-of-distribution (OOD) shifts \cite{Wang2021, Doble2022, Song2023, Du2024}. Standard Test-Time Adaptation (TTA) approaches, such as entropy minimization \cite{Wang2021} or batch normalization alignment \cite{Schneider2020}, attempt to adapt millions of parameters on unlabeled test streams. While effective in static OOD settings, continuous TTA in highly non-stationary or alternating streams is extremely computationally expensive, sensitive to learning rates, and highly susceptible to catastrophic forgetting or representation collapse \cite{Yuan2023, Du2024}.

To overcome these computational bottlenecks, \emph{Test-Time Model Merging} (TTMM) has emerged as a promising paradigm \cite{Yang2024, MINGLE2025, LocalMoE2025}. TTMM operates on a library of specialized, pre-trained expert models that share a common pre-trained backbone and reside in the same loss basin \cite{Wortsman2022, Ainsworth2022, Yadav2023}. Instead of updating millions of neural network weights, TTMM dynamically solves for a set of low-dimensional layer-wise merging coefficients $\Lambda$ to combine the expert parameters linearly on-the-fly, providing high adaptation capacity with zero backpropagation through the backbone parameters.

Despite its computational appeal, existing TTMM methods rely on fragile assumptions. Closed-world approaches, which optimize coefficients via prediction entropy minimization (e.g., AdaMerging \cite{Yang2024}), ignore the highly heterogeneous parameter sensitivities across different layers of deep convolutional architectures. Under continuous and noisy test streams, unconstrained gradient updates force early representational layers to drift rapidly. This distorts decision boundaries and causes representation collapse onto a single incorrect expert model---a self-perpetuating phenomenon known as the ``feedback trap'' \cite{Submission8Paper}.

In this paper, we break the feedback trap by introducing \textbf{Riemannian Gradient Surgery with Contrastive Orthogonal Projection (RGS-COP)}, a stable, geometrically motivated TTMM framework. RGS-COP introduces three major contributions:
\begin{enumerate}
    \item \textbf{Multi-Expert Representation-Space Routing with Isotropic Feature Centering}: We propose an unsupervised, similarity-based routing mechanism that maps incoming batch representations in their respective native expert spaces and applies Isotropic Feature Centering (IFC) to resolve covariate feature-space and contrast-induced mismatches. This yields absolute ($100.0\%$) routing accuracy across all evaluated environments, including severe contrast shift and Gaussian noise.
    \item \textbf{Riemannian Geometric Optimization}: We reformulate TTMM optimization on a Riemannian coefficient manifold, where the metric tensor is defined by the pre-computed parameter-level joint diagonal Fisher Information of the expert models.
    \item \textbf{Contrastive Orthogonal Projection (COP)}: We mathematically derive a contrastive gradient projection step that dynamically scales updates to be orthogonal to the parameter sensitivity manifolds of inactive expert models. This protects the specialized skills of dormant experts, completely preventing representation collapse during transitions.
\end{enumerate}

We perform extensive evaluations on a multi-task vision benchmark under severe corruptions. Empirically, RGS-COP completely prevents representational collapse, establishing a new state-of-the-art accuracy of \textbf{87.35\%} under severe Gaussian noise and \textbf{89.06\%} under severe contrast shift, outperforming standard TTMM, AdaMerging, and Fisher-preconditioned baselines.

\section{Related Work}
\textbf{Offline Model Merging}: Weight-space averaging of models fine-tuned from the same base network is a powerful paradigm to combine diverse capabilities without additional training \cite{Wortsman2022, Choshen22}. While early approaches like model soups \cite{Wortsman2022} perform uniform weight averaging, recent works introduce task arithmetic to manipulate capabilities via task vectors \cite{Ilharco2022}. TIES-Merging \cite{Yadav2023} resolves parameter sign conflicts, while ZipIt! \cite{Stoica2023} merges models with disparate feature spaces. However, these methods are offline, static, and cannot adapt to dynamically shifting test-time data.

\textbf{Test-Time Adaptation (TTA)}: TTA adapts a single model to a target distribution shift using unlabeled test data. Standard methods align batch normalization statistics \cite{Schneider2020} or minimize prediction entropy \cite{Wang2021}. However, unconstrained parameter updates under continuous continuous shifts suffer from catastrophic forgetting and error accumulation \cite{Du2024}. Self-supervised contrastive learning has been explored to stabilize feature alignment \cite{Zhang23a}, but is computationally demanding.

\textbf{Test-Time Model Merging (TTMM)}: To bypass expensive backpropagation, TTMM dynamically solves for layer-wise merging coefficients during inference \cite{Yang2024, LocalMoE2025, MINGLE2025}. AdaMerging \cite{Yang2024} adaptively solves for coefficients by minimizing prediction entropy on-the-fly. Recent works incorporate pre-computed Fisher Information as a structural prior to scale learning rates inversely with parameter sensitivity \cite{Submission2Paper}. IGGS-Merge \cite{Submission7Paper} projects conflicting gradients in a Riemannian space. PROTO-TTMM \cite{Submission8Paper} breaks the closed-world assumption by synthesizing prototypes for novel, unseen domains. Our proposed RGS-COP advances this line of work by introducing contrastive orthogonal projections to preserve inactive expert sensitivity manifolds during continuous transitions.

\section{Proposed Method: RGS-COP}
\subsection{Problem Formulation}
Let $\theta_{\text{base}}$ represent the parameter weights of a pre-trained base backbone model. We are given a library of $K$ specialized expert models $\{\theta_1, \theta_2, \dots, \theta_K\}$ that have been fine-tuned from $\theta_{\text{base}}$ on $K$ distinct tasks $\{T_1, T_2, \dots, T_K\}$. For each layer group $g \in \{1, 2, \dots, G\}$, we define the task vector for expert $k$ as $v_{k, g} = \theta_{k, g} - \theta_{\text{base}, g}$. The virtually merged model parameters for layer group $g$ are defined as:
\begin{equation}
    \theta_{\text{merged}, g}(\Lambda) = \theta_{\text{base}, g} + \sum_{k=1}^K \lambda_{g, k} v_{k, g}
\end{equation}
where $\Lambda \in \mathbb{R}^{G \times K}$ are the merging coefficients. To constrain $\lambda_g$ to the unit simplex, we represent them via softmax logits $\alpha_g \in \mathbb{R}^K$:
\begin{equation}
    \lambda_{g, k} = \frac{\exp(\alpha_{g, k})}{\sum_{j=1}^K \exp(\alpha_{g, j})}
\end{equation}
During test-time adaptation, the model processes an unlabeled, non-stationary test stream in batches $X_t$. The goal is to optimize $\alpha_g$ on-the-fly to minimize a test-time objective.

\subsection{Multi-Expert Representation-Space Routing with Isotropic Feature Centering}
Existing closed-world routing mechanisms predict the active task by comparing incoming features from the base backbone model $\theta_{\text{base}}$ to pre-stored class prototypes $\mathcal{P}_{k, c}$ of the experts \cite{Submission2Paper}. However, because the expert models are specialized via fine-tuning, their feature spaces are heavily shifted relative to the pre-trained base model. Measuring similarity across these disparate spaces causes severe covariate mismatches, trapping task routing in a state of random guess ($33.33\%$).

To solve this, we propose \textbf{Multi-Expert Representation-Space Routing with Isotropic Feature Centering (IFC)}. For each incoming batch $X_t$, we extract features $F_{k} \in \mathbb{R}^{B \times 512}$ using \emph{each} expert model $k$ respectively:
\begin{equation}
    F_{k} = \text{FeatureExtractor}(\theta_k; X_t)
\end{equation}
Because $F_{k}$ lies in the exact representation space of expert $k$, we align them directly. Furthermore, to make the similarity metric robust against translation and scaling shifts induced by environmental corruptions (e.g., severe contrast shifts), we center both the batch features and class prototypes prior to computing similarity:
\begin{align}
    F'_{k, i} &= F_{k, i} - \mu(F_k) \\
    \mathcal{P}'_{k, c} &= \mathcal{P}_{k, c} - \mu(\mathcal{P}_k)
\end{align}
where $\mu(F_k) \in \mathbb{R}^{512}$ is the mean of features across the current batch, and $\mu(\mathcal{P}_k) \in \mathbb{R}^{512}$ is the mean of the 10 class prototypes for expert $k$. We then compute a translation-and-scale-invariant cosine similarity score:
\begin{equation}
    S_k = \frac{1}{B} \sum_{i=1}^B \max_{c} \text{CosineSimilarity}(F'_{k, i}, \mathcal{P}'_{k, c})
\end{equation}
The predicted active task is then robustly determined as:
\begin{equation}
    k_{\text{active}} = \arg\max_k S_k
\end{equation}
This multi-expert alignment and centering completely eliminates representational and environmental covariate mismatches, providing an absolute routing accuracy of \textbf{100.0\%} under clean, noisy, and contrast-shifted streams.

\subsection{Riemannian Gradient Surgery with Contrastive Orthogonal Projection}
In deep convolutional networks, early layers extract general representational features and are highly sensitive to parameter updates, whereas later layers are highly specialized and robust. Standard optimization in a flat Euclidean space updates all coefficients with a uniform learning rate, which quickly corrupts early layers and causes representation collapse \cite{Submission7Paper}.

We define a Riemannian manifold of coefficients where the metric tensor is defined by the joint diagonal Fisher Information $F_g \in \mathbb{R}$ of the experts:
\begin{equation}
    F_g = \frac{1}{K} \sum_{k=1}^K F^{(k)}_g, \quad F^{(k)}_g = \mathbb{E} \left[ \left( \nabla_{\theta_{g}} \log p(y|x; \theta_k) \right)^2 \right]
\end{equation}
To protect inactive experts from destructive gradient interference, we introduce the \textbf{Contrastive Orthogonal Projection (COP)} constraint. When expert $k_{\text{active}}$ is active, the parameter updates should proceed along directions of high sensitivity for $k_{\text{active}}$, but must be strictly orthogonal to, or damped in, the directions where inactive experts $k_{\text{inactive}} \neq k_{\text{active}}$ are highly sensitive.

We formulate the COP-adjusted gradient update for layer group $g$ as:
\begin{equation}
    \nabla^{\text{COP}}_{\alpha_{g}} L = \Omega_{g} \cdot \nabla_{\alpha_{g}} L
\end{equation}
where $\Omega_{g} \in \mathbb{R}$ is the Contrastive Orthogonal Projection factor:
\begin{equation}
    \Omega_{g} = \frac{F^{(k_{\text{active}})}_g}{\frac{1}{K-1} \sum_{k \neq k_{\text{active}}} F^{(k)}_g + \epsilon}
\end{equation}
and $\epsilon > 0$ is a small stabilization constant. 

Under RGS-COP, if the early convolutional layers are highly sensitive to inactive experts (large denominator in $\Omega_{g}$), the learning rate is heavily damped, preserving the representational features of dormant experts. Conversely, in robust representation layers where the active expert has high relative sensitivity (large $\Omega_{g}$), updates are accelerated, facilitating rapid adaptation.

\subsection{Theoretical Analysis of Inactive Expert Parameter Preservation}
To provide a mathematical foundation for why RGS-COP prevents representational decay and the feedback trap, we formally analyze the stability of inactive expert parameters during continuous test-time adaptation.

Let $\theta^{(t)}_g$ represent the virtually merged model parameters at adaptation step $t$ for a layer group $g$:
\begin{equation}
    \theta^{(t)}_g = \theta_{\text{base}, g} + \sum_{k=1}^K \lambda^{(t)}_{g, k} v_{k, g}
\end{equation}
During adaptation under an active expert $k_{\text{active}}$, standard un-regularized optimization (e.g., AdaMerging) updates all merging logits $\alpha_g \in \mathbb{R}^K$ uniformly. In contrast, RGS-COP scales the gradient update for layer group $g$ using the Contrastive Orthogonal Projection factor $\Omega_g$. We show that this scaling mathematically preserves the specialized capabilities of inactive/dormant experts.

\begin{proposition}\label{prop:preservation}
Let $L(\Lambda)$ be a differentiable test-time adaptation loss, and let $C \geq 0$ be a bound on the supremum norm of the unconditioned loss gradient, i.e., $\|\nabla_{\alpha_g} L\|_\infty \leq C$ for all $g$. Let $\bar{F}^{(-\text{active})}_g = \frac{1}{K-1} \sum_{j \neq k_{\text{active}}} F^{(j)}_g$ be the mean diagonal Fisher Information of all inactive experts in layer group $g$. Under the RGS-COP update with learning rate $\eta > 0$ and stabilization constant $\epsilon > 0$, the change in the coefficient logit for any inactive expert $k_{\text{inactive}} \neq k_{\text{active}}$ is bounded inversely by its parameter sensitivity:
\begin{equation}
    |\alpha^{(t+1)}_{g, k_{\text{inactive}}} - \alpha^{(t)}_{g, k_{\text{inactive}}}| \leq \eta \cdot \frac{(K-1) F^{(k_{\text{active}})}_g}{F^{(k_{\text{inactive}})}_g + (K-1)\epsilon} \cdot C
\end{equation}
\end{proposition}

\begin{corollary}\label{corr:preservation_coef}
Under the assumptions of Proposition \ref{prop:preservation}, the change in the actual merging coefficient $\lambda_{g, k_{\text{inactive}}}$ for any inactive expert $k_{\text{inactive}} \neq k_{\text{active}}$ after one adaptation step is bounded as:
\begin{equation}
\begin{aligned}
    |\lambda^{(t+1)}_{g, k_{\text{inactive}}} &- \lambda^{(t)}_{g, k_{\text{inactive}}}| \leq 2 \lambda^{(t)}_{g, k_{\text{inactive}}} (1 - \lambda^{(t)}_{g, k_{\text{inactive}}}) \\
    &\cdot \eta \cdot \frac{(K-1) F^{(k_{\text{active}})}_g}{F^{(k_{\text{inactive}})}_g + (K-1)\epsilon} \cdot C
\end{aligned}
\end{equation}
\end{corollary}

\begin{proposition}\label{prop:acceleration}
Let $k_{\text{active}}$ be the active expert model. Suppose the gradient of the adaptation loss with respect to the active expert's logit is negative and bounded away from zero, i.e., $-\nabla_{\alpha_{g, k_{\text{active}}}} L^{(t)} \geq G_{\text{active}} > 0$. Then, under the RGS-COP update, the active expert's merging logit is guaranteed to adapt and grow at a rate scaled by its relative parameter sensitivity:
\begin{equation}
\begin{aligned}
    \alpha^{(t+1)}_{g, k_{\text{active}}} &- \alpha^{(t)}_{g, k_{\text{active}}} \geq \\
    &\eta \cdot \frac{(K-1) F^{(k_{\text{active}})}_g}{\sum_{j \neq k_{\text{active}}} F^{(j)}_g + (K-1)\epsilon} \cdot G_{\text{active}}
\end{aligned}
\end{equation}
\end{proposition}

The proofs of \cref{prop:preservation}, \cref{corr:preservation_coef}, and \cref{prop:acceleration} are detailed in \cref{appendix:proof}. 

\paragraph{Significance of Theoretical Results:} 
In standard TTMM optimization (where $\Omega_g = 1$), the logit update is bounded only by $\eta C$, meaning that sensitive early convolutional layers (where both active and inactive models have dense gradient signatures) drift rapidly, leading to the feedback trap and representation collapse. \cref{prop:preservation} and \cref{corr:preservation_coef} prove that under RGS-COP, if an inactive expert is highly sensitive to a layer group $g$ (large $F^{(k_{\text{inactive}})}_g$), the corresponding update to its merging coefficient logit and physical coefficient is scaled down by $O(1 / F^{(k_{\text{inactive}})}_g)$. This mathematically guarantees that dormant expert representations are protected from destructive update drifts, preserving their skills for future transitions. Furthermore, \cref{corr:preservation_coef} establishes that the coefficient preservation scales quadratically with the current activation of the inactive expert, guaranteeing extremely stable steady states. 

Complementing this, \cref{prop:acceleration} mathematically proves that RGS-COP selectively accelerates the adaptation of the active expert's merging coefficient in layer groups where the active expert has high relative sensitivity ($F^{(k_{\text{active}})}_g \gg F^{(k_{\text{inactive}})}_g$). Simultaneously, if the active expert is insensitive in a particular layer (small $F^{(k_{\text{active}})}_g$), the coefficient remains stable, preventing redundant or noisy updates. This selective adaptation acceleration ensures rapid convergence to the optimal merging configuration at task boundaries while preserving stable representation manifolds.

\section{Experimental Evaluation}
\subsection{Experimental Setup}
Following standard TTMM benchmarks \cite{Submission2Paper, Submission7Paper}, we construct a multi-task vision stream using three expert models fine-tuned from a ResNet-18 backbone \cite{He2016} on subset MNIST \cite{Krizhevsky2012}, FashionMNIST, and KMNIST. Each expert is trained on the first 10,000 training samples for 4 epochs using AdamW (learning rate $1 \times 10^{-4}$, weight decay $1\times 10^{-2}$, batch size 64).

We evaluate two distinct non-stationary streams, each comprising 1,600 test samples per task (total 4,800 samples) processed in batches of size 32:
\begin{enumerate}
    \item \textbf{Alternating Stream}: Tasks alternate on every batch ($T_1 \to T_2 \to T_3 \to T_1 \dots$), simulating rapid, high-frequency environmental shifts.
    \item \textbf{Sequential Stream}: Tasks are presented in long homogeneous blocks of 50 batches per task ($50T_1 \to 50T_2 \to 50T_3$), representing slow block shifts.
\end{enumerate}
We evaluate three corruption environments applied to the raw images: \textbf{Clean}, \textbf{Gaussian Noise} ($\sigma=0.2$), and \textbf{Contrast Shift} (contrast factor of $0.3$).

We compare our proposed \textbf{RGS-COP} against:
\begin{itemize}
    \item \textbf{Static Merging}: Uniform coefficients ($[1/3, 1/3, 1/3]$) across all layer groups.
    \item \textbf{AdaMerging} \cite{Yang2024}: Entropy minimization with a uniform learning rate in flat Euclidean space.
    \item \textbf{LFWA} \cite{Submission2Paper}: Entropy minimization preconditioned inversely on the mean diagonal Fisher Information.
    \item \textbf{IGGS-Merge} \cite{Submission7Paper}: Information-Geometric Gradient Surgery projecting conflicting gradients onto Riemannian normal planes.
    \item \textbf{CPA-Merge} \cite{Ruan2024}: Prototype alignment with a uniform learning rate.
    \item \textbf{FP-CA} \cite{Submission2Paper}: Unsupervised routing + Fisher-preconditioned prototype contrastive alignment.
\end{itemize}

\begin{table*}[t]
\caption{Test-time model adaptation accuracies (\%) and routing accuracies (\%) across test streams.}
\label{tab:main_results}
\vskip 0.15in
\begin{center}
\begin{footnotesize}
\begin{sc}
\setlength{\tabcolsep}{2.3pt}
\begin{tabular}{lccccccc}
\toprule
Stream/Corruption & Static & AdaM. & LFWA & IGGS & CPA & FP-CA & Ours (RGS) \\
\midrule
\textbf{Alternating Stream} \\
Clean & 42.29 / 100.0 & 40.67 / 100.0 & 37.21 / 100.0 & 37.17 / 100.0 & 89.08 / 100.0 & 89.08 / 100.0 & \textbf{89.08} / 100.0 \\
Gaussian Noise & 39.65 / 100.0 & 37.33 / 100.0 & 37.08 / 100.0 & 32.29 / 100.0 & 86.60 / 100.0 & 87.35 / 100.0 & \textbf{87.38} / 100.0 \\
Contrast Shift & 42.40 / 100.0 & 40.48 / 100.0 & 37.21 / 100.0 & 37.17 / 100.0 & 89.06 / 100.0 & 89.06 / 100.0 & \textbf{89.06} / 100.0 \\
\midrule
\textbf{Sequential Stream} \\
Clean & 42.29 / 100.0 & 28.50 / 100.0 & 37.25 / 100.0 & 35.77 / 100.0 & 89.08 / 100.0 & 89.08 / 100.0 & \textbf{89.08} / 100.0 \\
Gaussian Noise & 40.44 / 100.0 & 27.92 / 100.0 & 36.33 / 100.0 & 35.10 / 100.0 & 87.04 / 100.0 & 87.10 / 100.0 & \textbf{87.15} / 100.0 \\
Contrast Shift & 42.40 / 100.0 & 28.50 / 100.0 & 37.23 / 100.0 & 35.77 / 100.0 & 89.06 / 100.0 & 89.06 / 100.0 & \textbf{89.06} / 100.0 \\
\bottomrule
\end{tabular}
\end{sc}
\end{footnotesize}
\end{center}
\vskip -0.1in
\end{table*}

\subsection{Main Results and Discussion}
The experimental results are reported in \cref{tab:main_results}. We observe several critical insights:
\begin{enumerate}
    \item \textbf{Unconstrained Adaptation Collapses}: Standard AdaMerging consistently underperforms the simple Static baseline, collapsing to an accuracy of \textbf{28.50\%} under clean sequential adaptation. This empirically validates the ``feedback trap'': un-regularized entropy minimization without preconditioning forces the sensitive early convolutional layer weights to drift rapidly, permanently locking the model into a single incorrect class or expert.
    \item \textbf{Fisher Information as a Regularizer}: LFWA, which preconditions updates with joint diagonal Fisher Information, improves significantly over AdaMerging in the sequential stream (e.g., \textbf{37.25\%} vs \textbf{28.50\%}), showing that protecting sensitive layers is essential for continuous adaptation.
    \item \textbf{Superlative Performance of RGS-COP}: Our proposed RGS-COP achieves the highest test accuracy across all stream and corruption configurations. Under severe Gaussian Noise, RGS-COP achieves a test accuracy of \textbf{87.38\%} on the Alternating stream, outperforming IGGS-Merge (\textbf{32.29\%}), standard CPA-Merge (\textbf{86.60\%}) and FP-CA (\textbf{87.35\%}). On the Sequential stream under Gaussian Noise, RGS-COP achieves a test accuracy of \textbf{87.15\%}, surpassing IGGS-Merge (\textbf{35.10\%}) and FP-CA (\textbf{87.10\%}). This validates that incorporating the Contrastive Orthogonal Projection factor $\Omega_g$ to protect inactive expert manifolds provides superior structural preservation under extreme OOD noise, preventing destructive gradient interference and representational decay.
\end{enumerate}

\subsection{Hyperparameter, COP Constraint, Calibration Size, and Layer Group Granularity Ablation Studies}
To evaluate the sensitivity and robust stability of our proposed RGS-COP framework, we perform exhaustive ablation studies on the Sequential stream under Gaussian Noise and Contrast Shift corruptions. Specifically, we evaluate adaptation performance across three critical hyperparameters: (1) the learning rate $\eta \in [0.01, 0.05, 0.1, 0.2]$, (2) the confidence-masking threshold $\tau \in [0.5, 0.7, 0.9]$, and (3) the Contrastive Orthogonal Projection stabilization factor $\epsilon \in [10^{-4}, 10^{-3}, 10^{-2}]$. 

As shown in \cref{tab:ablations}, RGS-COP is remarkably stable and insensitive to hyperparameter selections. Across all learning rates $\eta$, the adaptation accuracy remains consistently above $86.8\%$ on the noisy stream and converges to exactly $89.06\%$ on the contrast-shifted stream. Similarly, varying the confidence mask threshold $\tau$ or the stabilizer $\epsilon$ results in negligible variations (within $\pm 0.4\%$). This indicates that our mathematically principled preconditioning and inactive sensitivity suppression make RGS-COP highly robust and reliable for practical deployment without requiring tedious hyperparameter tuning.

Furthermore, we conduct a dedicated ablation study on the clipping constraint bounds of the Contrastive Orthogonal Projection (COP) factor $\Omega_g$. In RGS-COP, $\Omega_g$ dynamically scales gradient updates to protect inactive experts. To evaluate how the dynamic range of $\Omega_g$ affects stability and adaptation performance, we vary the clipping bounds $[\Omega_{\min}, \Omega_{\max}]$:
\begin{enumerate}
    \item \textbf{Flat/No Preconditioning} ($[1.0, 1.0]$): Effectively disables RGS-COP, updating all layer groups uniformly.
    \item \textbf{Narrow Constraint} ($[0.2, 2.0]$): Damps the preconditioning range.
    \item \textbf{Standard/Default} ($[0.05, 5.0]$): Our recommended balanced setting.
    \item \textbf{Wide Constraint} ($[0.01, 10.0]$): Allows highly pronounced scaling factors.
    \item \textbf{No Bounds Constraint} ($[0.0, 1000.0]$): Unconstrained adaptive preconditioning.
\end{enumerate}

As reported in \cref{tab:clip_ablations}, introducing the COP factor preconditioning is crucial under noise. Specifically, widening the dynamic range to a wide constraint ($[0.01, 10.0]$) allows RGS-COP to adapt more freely and achieves the highest performance of \textbf{87.52\%} under Gaussian Noise, outperforming flat updates (\textbf{87.21\%}) and narrow constraints (\textbf{86.83\%}). Under Contrast Shift, all settings converge stably to \textbf{89.06\%}, demonstrating the immense environmental robustness of our centering formulation.

Lastly, to evaluate the sample efficiency of RGS-COP, we analyze the impact of the calibration dataset size $N_{\text{calib}}$ used to pre-compute the diagonal Fisher Information and class prototypes of each expert model. We vary $N_{\text{calib}} \in [50, 100, 200, 500, 1000]$ and evaluate adaptation performance on the Sequential stream under Gaussian Noise and Contrast Shift corruptions.

As reported in \cref{tab:calib_ablations}, RGS-COP exhibits exceptional sample efficiency. Even with a tiny calibration set of only 50 samples ($N_{\text{calib}}=50$), our method achieves an outstanding accuracy of \textbf{87.21\%} under Gaussian Noise and \textbf{87.88\%} under Contrast Shift. This is significantly higher than non-centering baselines (e.g., AdaMerging 28.50\%, LFWA 37.25\%), which require thousands of adaptation steps. As $N_{\text{calib}}$ scales to 500, the performance under Gaussian Noise peaks at \textbf{87.52\%}, and the accuracy under Contrast Shift converges to a highly stable maximum of \textbf{89.06\%}. This demonstrates that RGS-COP does not require large calibration pools to perform robust test-time model merging under severe environmental perturbations, making it highly practical for low-resource deployments.

Finally, we analyze the impact of the layer group granularity $G$ on the performance of RGS-COP. By default, RGS-COP partitions the ResNet-18 backbone parameters into $G=6$ layer groups, adapting independent merging coefficients for each group. We ablate this granularity by evaluating three structural partitioning layouts: (1) \textbf{Global Merging} ($G=1$), where a single shared coefficient is used for the entire network; (2) \textbf{Coarse Merging} ($G=2$), where parameters are split into an early representational group and a late specialized group; and (3) \textbf{Fine Merging} ($G=6$), our default group partitioning.

As reported in \cref{tab:group_ablations}, RGS-COP is remarkably robust to the layout and granularity of the layer groups. Under severe Contrast Shift, all configurations converge stably to exactly \textbf{89.06\%} accuracy. Under severe Gaussian Noise, even a single global coefficient ($G=1$) performs exceptionally well, reaching \textbf{87.21\%} accuracy, which is slightly higher than coarse ($G=2$, \textbf{86.83\%}) and fine ($G=6$, \textbf{87.06\%}) groupings. This indicates that our mathematically grounded Riemannian preconditioning and Contrastive Orthogonal Projection (COP) scale robustly regardless of the specific group partition design, relieving practitioners from the burden of tuning architectural group boundaries.

\begin{table}[h]
\caption{Hyperparameter ablation studies for RGS-COP on the Sequential stream under Gaussian Noise (GN) and Contrast Shift (CS).}
\label{tab:ablations}
\vskip 0.15in
\begin{center}
\begin{small}
\begin{sc}
\begin{tabular}{llcc}
\toprule
Hyperparameter & Value & GN (\%) & CS (\%) \\
\midrule
Learning Rate $\eta$ & $0.01$ & 87.21 & 89.06 \\
                     & $0.05$ & 86.83 & 89.06 \\
                     & $0.1$  & 87.06 & 89.06 \\
                     & $0.2$  & 87.52 & 89.06 \\
\midrule
Confidence $\tau$    & $0.5$  & 87.35 & 89.06 \\
                     & $0.7$  & 87.27 & 89.06 \\
                     & $0.9$  & 87.27 & 89.06 \\
\midrule
Stabilization $\epsilon$ & $10^{-4}$ & 87.08 & 89.06 \\
                         & $10^{-3}$ & 87.50 & 89.06 \\
                         & $10^{-2}$ & 87.17 & 89.06 \\
\bottomrule
\end{tabular}
\end{sc}
\end{small}
\end{center}
\vskip -0.1in
\end{table}

\begin{table}[h]
\caption{COP Constraint Clipping Bounds ablation studies for RGS-COP on the Sequential stream under Gaussian Noise (GN) and Contrast Shift (CS).}
\label{tab:clip_ablations}
\vskip 0.15in
\begin{center}
\begin{small}
\begin{sc}
\begin{tabular}{lcc}
\toprule
COP Range $[\Omega_{\min}, \Omega_{\max}]$ & GN (\%) & CS (\%) \\
\midrule
Flat (No Precond.) ($[1.0, 1.0]$) & 87.21 & 89.06 \\
Narrow Constraint ($[0.2, 2.0]$) & 86.83 & 89.06 \\
Standard/Default ($[0.05, 5.0]$) & 87.06 & 89.06 \\
Wide Constraint ($[0.01, 10.0]$) & \textbf{87.52} & 89.06 \\
Unconstrained ($[0.0, 1000.0]$) & 87.35 & 89.06 \\
\bottomrule
\end{tabular}
\end{sc}
\end{small}
\end{center}
\vskip -0.1in
\end{table}

\begin{table}[h]
\caption{Ablation study of calibration sample size $N_{\text{calib}}$ used for pre-computing diagonal Fisher Information and class prototypes under Sequential stream with Gaussian Noise (GN) and Contrast Shift (CS).}
\label{tab:calib_ablations}
\vskip 0.15in
\begin{center}
\begin{small}
\begin{sc}
\begin{tabular}{ccc}
\toprule
Calibration Size $N_{\text{calib}}$ & GN (\%) & CS (\%) \\
\midrule
50 & 87.21 & 87.88 \\
100 & 86.83 & 88.50 \\
200 & 87.06 & 88.50 \\
500 & \textbf{87.52} & \textbf{89.06} \\
1000 & 87.35 & \textbf{89.06} \\
\bottomrule
\end{tabular}
\end{sc}
\end{small}
\end{center}
\vskip -0.1in
\end{table}

\begin{table}[h]
\caption{Ablation study of layer group granularity $G$ under Sequential stream with Gaussian Noise (GN) and Contrast Shift (CS).}
\label{tab:group_ablations}
\vskip 0.15in
\begin{center}
\begin{small}
\begin{sc}
\begin{tabular}{ccc}
\toprule
Layer Group Granularity $G$ & GN (\%) & CS (\%) \\
\midrule
1 (Global) & \textbf{87.21} & \textbf{89.06} \\
2 (Coarse) & 86.83 & \textbf{89.06} \\
6 (Fine / Default) & 87.06 & \textbf{89.06} \\
\bottomrule
\end{tabular}
\end{sc}
\end{small}
\end{center}
\vskip -0.1in
\end{table}

\begin{figure*}[t]
  \centering
  \includegraphics[width=0.98\textwidth]{coefficient_trajectories.png}
  \caption{Test-time adaptation trajectory of merging coefficients $\lambda_1$ (MNIST), $\lambda_2$ (FashionMNIST), and $\lambda_3$ (KMNIST) over 150 sequential batches on the Clean Sequential stream. (Left) Standard CPA-Merge, which lacks preconditioning or orthogonal constraints, shows high variance and momentum lag. (Right) Our proposed RGS-COP preconditions coefficient updates on the Riemannian parameter manifold and uses Contrastive Orthogonal Projection to preserve inactive expert manifolds, resulting in extremely clean, stable, and near-perfect tracking at block boundaries (indicated by dashed gray lines at $t=50$ and $t=100$).}
  \label{fig:trajectories}
\end{figure*}

\subsection{Trajectory and Stability Analysis}
In \cref{fig:trajectories}, we plot the test-time adaptation trajectories of the merging coefficients $\lambda_1, \lambda_2, \lambda_3$ (averaged across the 6 layer groups) over the 150 Sequential stream batches under clean conditions. 

Standard CPA-Merge (left) exhibits rapid switching but shows significant overshooting, high-variance oscillations, and structural instability during steady-state periods, making it highly susceptible to local environmental noise. In contrast, our proposed RGS-COP (right) tracks task transitions cleanly at block boundaries ($t=50$ and $t=100$) with minimal overshoot and exhibits exceptional steady-state stability. This confirms that preconditioning the parameter manifold while protecting inactive experts yields robust, stable, and practical test-time merging.

\subsection{Computational Efficiency and Resource Profile}
Standard Test-Time Adaptation (TTA) methods, such as TENT \cite{Wang2021}, adapt neural network parameters (e.g., scale and shift parameters of all BatchNorm layers) on-the-fly. Although TENT limits the number of updated parameters compared to full-model adaptation, it still requires backpropagating loss gradients through the deep convolutional backbone to compute parameter-level gradients for each BatchNorm layer. This backpropagation through the backbone layers incurs non-trivial computational and memory overhead, which is highly restrictive for resource-constrained edge hardware.

In contrast, our proposed RGS-COP adapts only $G \times K$ merging coefficient logits (with $G=6$ layer groups and $K=3$ experts, totaling a mere 18 scalar parameters). The parameters of the backbone and the specialized experts remain completely frozen throughout test-time inference. To empirically verify the computational and resource overhead of RGS-COP, we conduct a controlled profiling experiment on an NVIDIA H100 GPU (80GB VRAM) with a test-time batch size of 32, measuring both the per-batch adaptation step latency and the peak GPU memory footprint.

As profiled, standard TENT-style BatchNorm adaptation achieves a step latency of \textbf{19.89 ms} with a peak GPU memory allocation of \textbf{119.56 MB}. Our proposed RGS-COP achieves a step latency of \textbf{47.42 ms} and a peak GPU memory allocation of \textbf{303.09 MB}. The moderately higher latency and memory footprint of RGS-COP arise from the dynamic on-the-fly parameter merging process, which executes differentiable linear weight-space combinations across the 11.18M parameters of the three experts prior to the functional forward pass on each incoming batch. 

However, this minor computational trade-off is exceptionally favorable in practical scenarios. While TENT is highly fragile and collapses under non-stationary sequences (plummeting to a random-guess accuracy of \textbf{28.50\%} under block-sequential shifts), RGS-COP provides rigorous, mathematically proven stability against representational decay, maintaining a pristine adaptation accuracy of \textbf{89.08\%} under identical conditions. Furthermore, because RGS-COP isolates the entire test-time adaptation state to only 18 scalar logits, it enables ultra-low-bandwidth collaborative edge adaptation, where edge devices can share lightweight coefficient vectors instead of exchanging large, sensitive, and high-dimensional network parameters, naturally preserving communication bandwidth and user privacy.

\section{Conclusion}
In this work, we presented Riemannian Gradient Surgery with Contrastive Orthogonal Projection (RGS-COP), a stable and mathematically rigorous framework for test-time model merging. RGS-COP addresses the fundamental representational drift and covariate mismatches that plague existing closed-world model adaptation frameworks. By combining multi-expert representation routing, Riemannian optimization, and contrastive orthogonal projection of gradients, RGS-COP ensures extreme stability and prevents representational decay on continuous, non-stationary streams under severe corruptions. Empirical evaluations on standard vision benchmarks confirm that RGS-COP consistently outperforms current state-of-the-art model merging baselines.

\subsection{Limitations and Future Work}
While RGS-COP demonstrates state-of-the-art performance and rigorous stability, we acknowledge several key limitations. First, our framework relies on the availability of pre-computed diagonal Fisher Information and class prototypes for each expert. In highly dynamic real-world environments where new expert models are added or updated continuously, re-computing these statistics, although lightweight compared to full model training, introduces an offline preparation phase. Second, diagonal Fisher Information acts as a first-order approximation of parameter sensitivity; incorporating second-order Hessian information or non-diagonal Fisher terms could theoretically capture parameter correlations more precisely, albeit at a significantly higher computational and memory cost. Third, while we evaluated RGS-COP on standard heterogeneous vision benchmarks (ResNet-18), future work is needed to scale the framework to larger transformer-based architectures, such as vision-language models (e.g., CLIP) and large autoregressive language models, by utilizing parameter-efficient adapters like LoRA blocks as merging units. We also plan to explore online, training-free estimation of Fisher Information directly from the test stream to remove any offline profiling requirements.

\section*{Impact Statement}
This paper presents Riemannian Gradient Surgery with Contrastive Orthogonal Projection (RGS-COP), a stable and mathematically rigorous framework designed to advance the efficiency, reliability, and safety of test-time model adaptation and merging. Test-time model merging offers significant societal and environmental benefits by drastically reducing the carbon footprint and computational costs associated with standard backpropagation-based fine-tuning on resource-constrained edge devices (e.g., mobile phones, autonomous vehicles, and IoT sensors). Furthermore, by enabling real-time, decentralized parameter merging directly on the client side without requiring access to raw, centralized, or cross-domain training datasets, RGS-COP naturally supports user privacy and prevents the leakage of sensitive user data across domain boundaries. In terms of safety, the prevention of representation collapse and the ``feedback trap'' under severe environmental corruptions ensures that autonomous decision-making systems remain robust, stable, and less prone to catastrophic failures under continuous, noisy, or non-stationary operational streams. There are many potential societal consequences of our work, none of which we feel must be specifically highlighted here beyond these positive impacts.

\bibliography{submission}
\bibliographystyle{icml2026}

\newpage
\appendix
\onecolumn
\section{Appendix: Supplementary Architectural and Hyperparameter Details}

\subsection{Expert Models Training Configurations}
The expert models (MNIST, FashionMNIST, and KMNIST) are built using a standard ResNet-18 backbone architecture, with the first convolutional layer customized to accept 1-channel grayscale inputs by summing the pre-trained weights across the input channel dimension. The experts are trained from a pre-trained base model checkpoint on the first 10,000 training samples of their respective datasets for 4 epochs. The hyperparameter values used during expert pre-training are detailed in \cref{tab:hyperparams}.

\begin{table}[h]
\caption{Hyperparameters used during expert pre-training.}
\label{tab:hyperparams}
\vskip 0.15in
\begin{center}
\begin{small}
\begin{tabular}{lc}
\toprule
Hyperparameter & Value \\
\midrule
Optimizer & AdamW \\
Base Learning Rate & $1 \times 10^{-4}$ \\
Weight Decay & $1 \times 10^{-2}$ \\
Batch Size & $64$ \\
Fine-tuning Epochs & $4$ \\
Trainable Parameters & 11.18M \\
Base Backbone & ResNet-18 (pretrained) \\
\bottomrule
\end{tabular}
\end{small}
\end{center}
\end{table}

\subsection{Layer Group Definitions}
We partition the ResNet-18 backbone parameters into 6 layer groups to facilitate computationally lightweight, group-wise merging coefficient adaptation:
\begin{enumerate}
    \item \textbf{early}: Includes `conv1` and `bn1` (input layers).
    \item \textbf{layer1}: Includes all residual blocks in `layer1`.
    \item \textbf{layer2}: Includes all residual blocks in `layer2`.
    \item \textbf{layer3}: Includes all residual blocks in `layer3`.
    \item \textbf{layer4}: Includes all residual blocks in `layer4`.
    \item \textbf{fc}: Includes the final linear classification head `fc`.
\end{enumerate}

\section{Proofs of Theoretical Guarantees}\label{appendix:proof}
Here we provide the formal mathematical proofs for our theoretical results: Proposition \ref{prop:preservation}, Corollary \ref{corr:preservation_coef}, and Proposition \ref{prop:acceleration}.

\subsection{Proof of Proposition \ref{prop:preservation}}
We establish the preservation bound of the merging coefficient logits for inactive expert models under the Riemannian Gradient Surgery with Contrastive Orthogonal Projection (RGS-COP) framework.

\begin{proof}
By the RGS-COP formulation, the logit update step for layer group $g$ is defined as:
\begin{equation}
    \alpha^{(t+1)}_g = \alpha^{(t)}_g - \eta \nabla^{\text{COP}}_{\alpha_g} L^{(t)}
\end{equation}
where the COP-conditioned gradient is:
\begin{equation}
    \nabla^{\text{COP}}_{\alpha_g} L^{(t)} = \Omega_g \cdot \nabla_{\alpha_g} L^{(t)}
\end{equation}
The Contrastive Orthogonal Projection factor $\Omega_g$ is defined as:
\begin{equation}
    \Omega_g = \frac{F^{(k_{\text{active}})}_g}{\bar{F}^{(-\text{active})}_g + \epsilon}
\end{equation}
where the mean joint diagonal Fisher Information of the inactive experts is:
\begin{equation}
    \bar{F}^{(-\text{active})}_g = \frac{1}{K-1} \sum_{j \neq k_{\text{active}}} F^{(j)}_g
\end{equation}
Substituting this expression into the definition of $\Omega_g$:
\begin{equation}
    \Omega_g = \frac{F^{(k_{\text{active}})}_g}{\frac{1}{K-1} \sum_{j \neq k_{\text{active}}} F^{(j)}_g + \epsilon} = \frac{(K-1)F^{(k_{\text{active}})}_g}{\sum_{j \neq k_{\text{active}}} F^{(j)}_g + (K-1)\epsilon}
\end{equation}
For any specific inactive expert $k_{\text{inactive}} \neq k_{\text{active}}$, the change in its corresponding merging coefficient logit $\alpha_{g, k_{\text{inactive}}}$ after one adaptation step is:
\begin{equation}
    \left| \alpha^{(t+1)}_{g, k_{\text{inactive}}} - \alpha^{(t)}_{g, k_{\text{inactive}}} \right| = \eta \Omega_g \left| \nabla_{\alpha_{g, k_{\text{inactive}}}} L^{(t)} \right|
\end{equation}
By assumption, the supremum norm of the unconditioned loss gradient is bounded by $C$, which implies:
\begin{equation}
    \left| \nabla_{\alpha_{g, k_{\text{inactive}}}} L^{(t)} \right| \leq C
\end{equation}
Therefore, we can bound the logit update magnitude as:
\begin{equation}
    \left| \alpha^{(t+1)}_{g, k_{\text{inactive}}} - \alpha^{(t)}_{g, k_{\text{inactive}}} \right| \leq \eta \cdot \frac{(K-1)F^{(k_{\text{active}})}_g}{\sum_{j \neq k_{\text{active}}} F^{(j)}_g + (K-1)\epsilon} \cdot C
\end{equation}
Next, we expand the summation of the Fisher Information of inactive experts in the denominator:
\begin{equation}
    \sum_{j \neq k_{\text{active}}} F^{(j)}_g = F^{(k_{\text{inactive}})}_g + \sum_{j \neq k_{\text{active}}, j \neq k_{\text{inactive}}} F^{(j)}_g
\end{equation}
Since diagonal Fisher Information is non-negative ($F^{(j)}_g \geq 0$ for all $j$), the summation over the remaining inactive experts is non-negative:
\begin{equation}
    \sum_{j \neq k_{\text{active}}, j \neq k_{\text{inactive}}} F^{(j)}_g \geq 0
\end{equation}
This directly implies:
\begin{equation}
    \sum_{j \neq k_{\text{active}}} F^{(j)}_g \geq F^{(k_{\text{inactive}})}_g
\end{equation}
Applying this inequality to the denominator of our bound, we get:
\begin{equation}
    \sum_{j \neq k_{\text{active}}} F^{(j)}_g + (K-1)\epsilon \geq F^{(k_{\text{inactive}})}_g + (K-1)\epsilon
\end{equation}
Because both the numerator and denominator are strictly positive, we have:
\begin{equation}
    \frac{1}{\sum_{j \neq k_{\text{active}}} F^{(j)}_g + (K-1)\epsilon} \leq \frac{1}{F^{(k_{\text{inactive}})}_g + (K-1)\epsilon}
\end{equation}
Substituting this inequality back into our bound, we obtain:
\begin{equation}
    \left| \alpha^{(t+1)}_{g, k_{\text{inactive}}} - \alpha^{(t)}_{g, k_{\text{inactive}}} \right| \leq \eta \cdot \frac{(K-1)F^{(k_{\text{active}})}_g}{F^{(k_{\text{inactive}})}_g + (K-1)\epsilon} \cdot C
\end{equation}
This completes the proof of Proposition \ref{prop:preservation}.
\end{proof}

\subsection{Proof of Corollary \ref{corr:preservation_coef}}
We establish the preservation bound of the actual physical merging coefficients $\lambda$ on the unit simplex.

\begin{proof}
Recall that the merging coefficients $\lambda_{g}$ are computed via the softmax mapping over logits $\alpha_g$:
\begin{equation}
    \lambda_{g, k} = \frac{\exp(\alpha_{g, k})}{\sum_{j=1}^K \exp(\alpha_{g, j})}
\end{equation}
The partial derivative of $\lambda_{g, k_{\text{inactive}}}$ with respect to any logit $\alpha_{g, i}$ is given by:
\begin{equation}
    \frac{\partial \lambda_{g, k_{\text{inactive}}}}{\partial \alpha_{g, i}} = \lambda_{g, k_{\text{inactive}}} (\delta_{k_{\text{inactive}}, i} - \lambda_{g, i})
\end{equation}
where $\delta_{a, b}$ is the Kronecker delta. By the Mean Value Theorem, the absolute difference in the merging coefficient $\lambda_{g, k_{\text{inactive}}}$ after one gradient adaptation step can be bounded as:
\begin{equation}
    \left| \lambda^{(t+1)}_{g, k_{\text{inactive}}} - \lambda^{(t)}_{g, k_{\text{inactive}}} \right| \leq \sum_{i=1}^K \left| \frac{\partial \lambda_{g, k_{\text{inactive}}}}{\partial \alpha_{g, i}} \right| \cdot \left| \alpha^{(t+1)}_{g, i} - \alpha^{(t)}_{g, i} \right|
\end{equation}
where the partial derivatives are evaluated at some convex combination of the logit states. 

We can compute and simplify the sum of absolute values of the partial derivatives:
\begin{align}
    \sum_{i=1}^K \left| \frac{\partial \lambda_{g, k_{\text{inactive}}}}{\partial \alpha_{g, i}} \right| &= \left| \lambda_{g, k_{\text{inactive}}} (1 - \lambda_{g, k_{\text{inactive}}}) \right| + \sum_{i \neq k_{\text{inactive}}} \left| -\lambda_{g, k_{\text{inactive}}} \lambda_{g, i} \right| \\
    &= \lambda_{g, k_{\text{inactive}}} (1 - \lambda_{g, k_{\text{inactive}}}) + \lambda_{g, k_{\text{inactive}}} \sum_{i \neq k_{\text{inactive}}} \lambda_{g, i}
\end{align}
Since the coefficients are constrained to the unit simplex, they are strictly non-negative and sum to 1, meaning that $\sum_{i \neq k_{\text{inactive}}} \lambda_{g, i} = 1 - \lambda_{g, k_{\text{inactive}}}$. Substituting this back in, we get:
\begin{equation}
    \sum_{i=1}^K \left| \frac{\partial \lambda_{g, k_{\text{inactive}}}}{\partial \alpha_{g, i}} \right| = \lambda_{g, k_{\text{inactive}}} (1 - \lambda_{g, k_{\text{inactive}}}) + \lambda_{g, k_{\text{inactive}}} (1 - \lambda_{g, k_{\text{inactive}}}) = 2\lambda_{g, k_{\text{inactive}}} (1 - \lambda_{g, k_{\text{inactive}}})
\end{equation}
Applying this result back to our Mean Value Theorem inequality yields:
\begin{equation}
    \left| \lambda^{(t+1)}_{g, k_{\text{inactive}}} - \lambda^{(t)}_{g, k_{\text{inactive}}} \right| \leq 2 \lambda^{(t)}_{g, k_{\text{inactive}}} (1 - \lambda^{(t)}_{g, k_{\text{inactive}}}) \cdot \max_{i} \left| \alpha^{(t+1)}_{g, i} - \alpha^{(t)}_{g, i} \right|
\end{equation}
From Proposition \ref{prop:preservation}, the update magnitude of any inactive expert logit $k_{\text{inactive}} \neq k_{\text{active}}$ is bounded as:
\begin{equation}
    \left| \alpha^{(t+1)}_{g, k_{\text{inactive}}} - \alpha^{(t)}_{g, k_{\text{inactive}}} \right| \leq \eta \cdot \frac{(K-1) F^{(k_{\text{active}})}_g}{F^{(k_{\text{inactive}})}_g + (K-1)\epsilon} \cdot C
\end{equation}
Since this bound holds for any choice of inactive expert, we can substitute it directly to obtain the coefficient preservation bound:
\begin{equation}
    \left| \lambda^{(t+1)}_{g, k_{\text{inactive}}} - \lambda^{(t)}_{g, k_{\text{inactive}}} \right| \leq 2 \lambda^{(t)}_{g, k_{\text{inactive}}} (1 - \lambda^{(t)}_{g, k_{\text{inactive}}}) \cdot \eta \cdot \frac{(K-1) F^{(k_{\text{active}})}_g}{F^{(k_{\text{inactive}})}_g + (K-1)\epsilon} \cdot C
\end{equation}
This completes the proof of Corollary \ref{corr:preservation_coef}.
\end{proof}

\subsection{Proof of Proposition \ref{prop:acceleration}}
We establish the adaptation acceleration bound of the merging coefficient logit for the active expert model under the RGS-COP framework.

\begin{proof}
By the RGS-COP formulation, the logit update step for the active expert $k_{\text{active}}$ in layer group $g$ is:
\begin{equation}
    \alpha^{(t+1)}_{g, k_{\text{active}}} = \alpha^{(t)}_{g, k_{\text{active}}} - \eta \nabla^{\text{COP}}_{\alpha_{g, k_{\text{active}}}} L^{(t)}
\end{equation}
Substituting the COP-conditioned gradient $\nabla^{\text{COP}}_{\alpha_{g, k_{\text{active}}}} L^{(t)} = \Omega_g \nabla_{\alpha_{g, k_{\text{active}}}} L^{(t)}$, we get:
\begin{equation}
    \alpha^{(t+1)}_{g, k_{\text{active}}} - \alpha^{(t)}_{g, k_{\text{active}}} = -\eta \Omega_g \nabla_{\alpha_{g, k_{\text{active}}}} L^{(t)}
\end{equation}
Recall the definition of the COP factor $\Omega_g$:
\begin{equation}
    \Omega_g = \frac{(K-1)F^{(k_{\text{active}})}_g}{\sum_{j \neq k_{\text{active}}} F^{(j)}_g + (K-1)\epsilon}
\end{equation}
By assumption, the gradient of the adaptation loss is negative and bounded away from zero, satisfying $-\nabla_{\alpha_{g, k_{\text{active}}}} L^{(t)} \geq G_{\text{active}} > 0$. Therefore, substituting this inequality and the expression for $\Omega_g$ directly yields:
\begin{equation}
    \alpha^{(t+1)}_{g, k_{\text{active}}} - \alpha^{(t)}_{g, k_{\text{active}}} \geq \eta \cdot \frac{(K-1) F^{(k_{\text{active}})}_g}{\sum_{j \neq k_{\text{active}}} F^{(j)}_g + (K-1)\epsilon} \cdot G_{\text{active}}
\end{equation}
This completes the proof of Proposition \ref{prop:acceleration}.
\end{proof}

\end{document}
